\documentclass[12pt,a4paper]{article}
\usepackage[UTF8,heading=true]{ctex}
\usepackage{geometry}
\usepackage{amsmath,amssymb,amsthm}
\usepackage{graphicx}
\usepackage{booktabs}
\usepackage{longtable}
\usepackage{array}
\usepackage{enumitem}
\usepackage{hyperref}
\usepackage{fancyhdr}
\usepackage{xcolor}
\usepackage{listings}
\usepackage{setspace}  % 添加行距控制包

% 配置 ctex 章节格式为 "一、"
\ctexset{
  section = {
    name = {,、},
    number = \chinese{section},
    format = {\Large\bfseries\raggedright}
  }
}

\geometry{left=2.5cm,right=2.5cm,top=2.5cm,bottom=2.5cm}
\setlength{\headheight}{15pt}
\setstretch{1.5}  % 设置1.5倍行距（CNIPA推荐）

\pagestyle{fancy}
\fancyhf{}
\fancyhead[C]{发明专利技术交底书}
\fancyfoot[C]{\thepage}

\title{\textbf{发明专利技术交底书} \\[0.5em]
\Large 面向大模型投机解码的非对称推理微架构及单周期回滚引擎方法与装置}
\author{中国移动通信有限公司研究院}
\date{\today}

\begin{document}

\begin{center}
    {\Large\heiti 中国移动专利申请} \\[1em]
    {\Huge\heiti 技术交底书} \\[1em]
\end{center}

\noindent
\begin{tabular}{|p{3cm}|p{12cm}|}
\hline
公司编号 &  \\
\hline
发明名称 & 面向大模型投机解码的非对称推理微架构及单周期回滚引擎方法与装置 \\
\hline
申报单位 & 研究院 \\
\hline
申报类型 & 发明 \\
\hline
发明人 & 许达 \\
\hline
技术联系人 & 许达 (xudayj@chinamobile.com, +86-13521894156) \\
\hline
注意事项 & 1．技术联系人应为深入了解本申请提案技术方案的技术人员。 \\
         & 2．请按照集团公司提供的本技术交底书模板逐项填写。 \\
\hline
\end{tabular}

\vspace{1cm}

%==============================================================================
\section{发明名称}
%==============================================================================

面向大模型投机解码的非对称推理微架构及单周期回滚引擎方法与装置


%==============================================================================
\section{技术领域}
%==============================================================================

本发明涉及人工智能计算、处理器微体系结构、异构多核加速器、片上网络（NoC）、存储层次结构、指令集与硬件状态机控制等技术领域，尤其涉及一种面向大语言模型（LLM）推理中的投机解码（Speculative Decoding）场景，将投机生成、验证、接受/拒绝判定以及提交与回滚控制下沉为硬件原生闭环的非对称微架构及方法与装置。

\textbf{与量子计算的关联说明：}

本发明不依赖新材料与量子器件，属于成熟 CMOS 数字逻辑可实现的芯片微架构与软硬协同设计。本发明的核心在于针对投机解码的控制路径和回滚路径进行 Algorithm-to-Silicon（算法硅化）：将原本由软件框架完成的“候选缓存、概率比较、拒绝采样、回滚与提交”等流程固化为硬件状态机与数据通路，从而提升推理吞吐并稳定延迟尾部。

本发明可应用于云端推理 GPU/TPU/ASIC、端侧 NPU/SoC 以及可重构加速器的解码阶段加速，尤其适用于交互式生成、多会话并发服务与能效敏感场景。

\textbf{关联项目：}\underline{\hspace{8cm}}（待填写）

%==============================================================================
\section{现有技术的技术方案}
%==============================================================================

\subsection{量子计算网络安全背景}

大模型推理通常采用自回归解码（Autoregressive Decoding），每生成一个 Token 需要基于前缀状态继续计算，存在强串行性。为了提升单序列 Tokens/sec，业界提出投机解码：由草稿模型（Draft Model）一次性生成 $k$ 个候选 Token，再由目标模型/验证模型（Target/Verify Model）对该候选序列进行概率验证，并根据接受准则决定提交候选前缀。

在通用硬件上，投机解码的典型流水包括：草稿生成、候选缓存、目标验证前向、候选对齐、逐 Token 接受判定、提交或回滚。该流程涉及大量非规则控制与动态分支，通常由软件运行时与驱动层组织并通过多个 kernel/算子序列实现。

\subsection{后量子密码技术现状}

现有软硬件实现投机解码通常采用同构计算阵列（例如 GPU 的 SIMT 核心）同时承担两类工作：
\begin{enumerate}
    \item \textbf{草稿推演计算}：执行小模型前向并产生候选；
    \item \textbf{验证与控制}：执行大模型验证前向，并在软件层进行候选对齐、概率比较、拒绝采样、以及失败后的回滚清理。
\end{enumerate}

在现有实现中，验证失败后的回滚通常需要软件显式处理：撤销未提交的 KV Cache 增量写入、回退序列位置指针、清空候选缓存、恢复 RNG 状态并重新触发草稿推演。上述回滚与比较逻辑在通用计算单元上执行效率低，且会引入不可预测的长尾延迟。

%==============================================================================
\section{现有技术的缺点及本申请提案要解决的技术问题}
%==============================================================================

现有技术至少存在以下缺点：
\begin{enumerate}
    \item \textbf{软件调度开销大}：草稿与验证切换导致频繁的 kernel launch、同步与队列管理，吞噬投机解码的理论收益；
    \item \textbf{非规则控制流低效}：逐 Token 的接受/拒绝判定与拒绝采样具有强分支性，易造成 SIMD/SIMT 低利用率；
    \item \textbf{回滚代价高且抖动大}：验证失败时需要软件介入清理状态与内存元数据，导致延迟尾部变差并引入缓存/显存抖动；
    \item \textbf{能效不佳}：将比较、随机数、回滚等控制路径压在通用算力单元上，导致单位 Token 能耗上升。
\end{enumerate}

因此，本申请提案要解决的技术问题是：\textbf{提供一种面向投机解码的非对称微架构，使草稿推演与目标验证在硬件层并行协同，并通过硬件影子队列与硬件拒绝采样器完成接受判定；当发生拒绝时，回滚引擎能够在单周期完成流水线清空、候选截断、程序计数器与提交边界恢复，从而实现软件零参与的快速回滚与稳定低延迟。}

%==============================================================================
\section{本申请提案的技术方案的详细阐述}
%==============================================================================

\subsection{系统架构}

本发明提出一种面向投机解码的硬件原生非对称推理微架构，包括：
\begin{enumerate}
    \item \textbf{草稿核（Draft Core）}：用于执行草稿模型前向与候选生成，配置为较高主频与较多实例，用于覆盖草稿推演的前向延迟；
    \item \textbf{验证核（Heavy/Verify Core）}：用于执行目标模型验证前向，配置为较强矩阵计算阵列与带宽，用于并行验证候选窗口；
    \item \textbf{影子队列（Shadow FIFO）}：用于承载草稿核生成的候选 Token 及其元数据（序列 ID、位置、可选概率摘要、RNG 状态摘要等）；
    \item \textbf{硬件拒绝采样器（HW Rejection Sampler）}：对验证核输出的目标分布与候选 Token 进行硬件比较与接受/拒绝判定；
    \item \textbf{提交/回滚控制器（Commit \& Rollback Controller）}：维护提交边界（commit frontier）、投机窗口指针、Shadow FIFO 的读写指针、以及投机增量区（例如 KV Cache 增量块）的资源登记，并在拒绝事件触发时执行单周期回滚；
    \item \textbf{投机状态表（Speculation State Table, SST）}：记录每个序列的投机参数（窗口 $k$、配额、回滚计数、检查点标识等），并支持多会话隔离与背压。
\end{enumerate}

在该架构中，草稿核与验证核并行执行。草稿核仅将候选写入 Shadow FIFO 并登记投机增量区占用，不直接改变提交态；验证核对候选窗口计算目标概率；硬件拒绝采样器做出接受/拒绝判定；控制器根据判定结果进行提交或回滚，且关键路径不依赖软件参与。

\subsection{核心方法步骤}

\subsubsection{1. 分布式密钥生成 (KeyGen)}

对应投机解码的初始化与检查点建立步骤。控制器为每个序列生成并维护“提交边界检查点”，其内容至少包括：序列位置、已提交的 KV Cache 指针、Shadow FIFO 边界指针、RNG/PRNG 状态标识与资源配额。该检查点用于后续拒绝时的单周期回滚恢复。

\subsubsection{2. 消息哈希与映射}

对应候选 Token 的结构化封装与对齐映射步骤。草稿核对每个候选 Token 生成元数据条目并写入 Shadow FIFO。为降低对齐开销，可对候选（Token ID，位置，序列 ID）计算短哈希或标签，用于验证核输出与候选的快速关联与一致性检查，从而避免软件层索引重排。

\subsubsection{3. 分布式高斯采样}

对应硬件拒绝采样与随机数生成步骤。硬件拒绝采样器基于验证核输出概率分布，对候选执行接受判定。所需随机数由硬件 RNG/PRNG 提供；为可复现与审计，可使用检查点中的 RNG 状态标识对随机数流进行确定性复位。采样器逐 Token 推进，并在首个拒绝点产生拒绝事件，携带拒绝位置与序列 ID。

\subsubsection{4. 基于 NTT 的线性变换}

对应验证计算的并行化执行与数据通路优化步骤。验证核可对候选窗口进行批量前向计算，并通过片上互连将结果以流式方式送入硬件拒绝采样器。为避免控制路径成为瓶颈，本发明可将候选窗口的对齐、读取与结果聚合固化为固定功能单元与状态机，使验证计算与判定流水化。

\subsubsection{5. 基于 Beaver 三元组的安全范数验证}

对应硬件一致性约束与“精确提交”保证步骤。提交/回滚控制器维护提交边界与投机边界，保证以下不变量：未提交的投机增量区不得污染提交态；提交时仅通过指针推进实现状态转移；拒绝时通过指针回退与流水线清空恢复到检查点。该机制确保在多会话并发下仍能实现精确提交与精确回滚。

\subsubsection{6. 动态节点管理}

对应多会话调度、背压与资源隔离步骤。SST 为每个序列提供独立条目，Shadow FIFO 支持每序列或每组序列隔离。控制器在资源紧张时可动态下调投机窗口 $k$ 或对草稿核施加背压，防止候选堆积与内存拥塞，并对回滚/提交事件进行计数以供性能分析（不进入关键路径）。

%==============================================================================
\section{本申请提案的关键点和欲保护点}
%==============================================================================

本申请提案主要包含以下关键创新点和保护点：
\begin{enumerate}
    \item \textbf{面向投机解码的非对称微架构}：将草稿推演与目标验证映射到不同类型核心（Draft Core 与 Heavy/Verify Core），并定义两者的并行协同机制；
    \item \textbf{硬件影子队列（Shadow FIFO）与候选元数据格式}：用于承载候选 Token、序列标识、位置标识、可选概率摘要与 RNG 状态摘要，并与投机增量区资源登记关联；
    \item \textbf{硬件拒绝采样器}：在硬件层执行候选的接受/拒绝判定，形成验证输出到提交/回滚控制的闭环；
    \item \textbf{单周期回滚引擎}：在拒绝事件触发时，单周期完成流水线清空、Shadow FIFO 截断、PC/序列位置恢复、投机增量区指针回收与提交边界恢复；
    \item \textbf{精确提交机制}：在接受前缀可提交时，通过指针推进将投机态原子切换为提交态，避免显式内存复制与软件清理；
    \item \textbf{多会话隔离与背压控制}：通过投机状态表与配额机制，实现并发服务下的稳定延迟与资源可控。
\end{enumerate}

%==============================================================================
\section{与第三条中最接近的现有技术相比，本申请提案有何技术优点}
%==============================================================================

相比于在通用 GPU/CPU 上以软件方式实现投机解码的现有技术，本申请提案具有如下技术优点：
\begin{enumerate}
    \item \textbf{吞吐更高且更稳定}：控制路径与判定路径硬件化，减少 kernel launch 与同步导致的间隙，使理论加速更可兑现；
    \item \textbf{显著降低延迟尾部}：回滚为单周期硬件动作，避免软件清理导致的长尾抖动；
    \item \textbf{更高能效}：将非规则控制与比较从通用算力单元剥离，减少无效计算与访存；
    \item \textbf{更强并发能力}：Shadow FIFO、SST 与背压机制提供会话隔离与拥塞控制，适合推理服务的多序列并行。
\end{enumerate}

%==============================================================================
\section{其他有助于理解本申请提案的技术资料}
%==============================================================================

(1) Leviathan, Y., Kalman, M., Matias, Y. ``Fast Inference from Transformers via Speculative Decoding.'' arXiv, 2023. \\
(2) Hennessy, J., Patterson, D. ``Computer Architecture: A Quantitative Approach.''（分支预测、检查点、回滚与精确异常相关章节）\\
(3) Yeh, T.-Y., Patt, Y.-N. ``Two-Level Adaptive Training Branch Prediction.'' MICRO, 1991. \\
(4) 典型推理服务系统资料：vLLM/Serving 系统公开技术报告（PagedAttention、调度与 KV 管理等）。

%==============================================================================
\section{本申请提案的侵权证据可获得性}
%==============================================================================

本提案的侵权行为具有较高的可检测性：
\begin{enumerate}
    \item \textbf{可观测性能特征}：单周期回滚会显著改变拒绝事件下的延迟分布，可通过基准测试与统计分析推断；
    \item \textbf{软硬件接口痕迹}：若实现硬件影子队列、硬件采样与回滚控制，驱动/编译器通常会暴露配置项、寄存器或性能计数器；
    \item \textbf{架构资料可对比}：产品白皮书、专利、SDK 文档与微基准可揭示非对称草稿/验证核与回滚机制的存在；
    \item \textbf{可复现实验}：通过控制投机窗口 $k$、拒绝率与并发度，观察吞吐与功耗曲线，可验证其是否采用硬件原生判定与回滚。
\end{enumerate}


\end{document}
